\documentclass{article}
\usepackage{PRIMEarxiv} % Core package for the PrimeAI Template
\usepackage{amsmath, amssymb, amsthm, bm, dcolumn} % Add amsthm here for the proof environment
\usepackage[numbers,sort&compress]{natbib} % Natbib for citations
\usepackage{graphicx} % For high-quality images
\usepackage{multicol} % Multi-column figures/tables
\usepackage{hyperref} % Hyperlinks
\usepackage{listings} % Code listings
\usepackage{authblk} % For structured affiliations
% Define theorem style (optional)
\theoremstyle{plain}
\newtheorem{theorem}{Theorem}
\renewcommand{\qedsymbol}{}

% Custom commands for primes and qubit encoding
\newcommand{\primeQubit}[1]{|p_{#1}\rangle}
\newcommand{\primeGate}[1]{U_{p_{#1}}}

\usepackage{wrapfig}
\usepackage[pscoord]{eso-pic}
\usepackage[fulladjust]{marginnote}
\reversemarginpar

% Typesetting improvements without footnote patching
\usepackage[protrusion=true, expansion=true, tracking=false]{microtype}
\microtypecontext{spacing=nonfrench}

% Line numbers
\usepackage[right]{lineno}

% Text layout - adjust as needed
\raggedright
\setlength{\parindent}{0.5cm}
\textwidth 5.25in 
\textheight 8.75in

% Set double spacing
\usepackage{setspace} 
\doublespacing

% Adjust width for specific content
\usepackage{changepage}

% Adjust caption style
\usepackage[aboveskip=1pt,labelfont=bf,labelsep=period,singlelinecheck=off]{caption}
% Define custom claims environment
\newtheoremstyle{claimstyle}
  {3pt}   % Space above
  {3pt}   % Space below
  {\itshape}   % Body font
  {}      % Indent amount
  {\bfseries} % Theorem head font
  {.}     % Punctuation after theorem head
  { }     % Space after theorem head
  {\thmname{#1}\thmnumber{ #2}\thmnote{ (#3)}} % Theorem head spec

\theoremstyle{claimstyle}
\newtheorem{claim}{Claim}
% Remove brackets from references
\makeatletter
\renewcommand{\@biblabel}[1]{\quad#1.}
\makeatother

% Header, footer, and page numbers
\usepackage{lastpage,fancyhdr}
\pagestyle{fancy}
\fancyhf{}
\fancyhead[L]{Preprint - PrimeAI Enhanced Template}
\fancyfoot[C]{\scriptsize Multiplicity Theory © 2024 Citizen Gardens \\ Licensed Under MIT and CC BY-NC-SA 4.0.}
\fancyfoot[R]{Page \thepage\ of \pageref{LastPage}}
\renewcommand{\footrule}{\hrule height 2pt \vspace{2mm}}


\title{Quantum Search Engine Development}
\author{C.R. Kunferman\\ \\ A Community Research Initiative\\ \large Citizen Gardens - \textit{The Foundation of Multiplicity}}
\date{\today}

\begin{document}
\maketitle

\section{Introduction}
This document presents the mathematical foundation for the integration of Kunferman's node-based parsing system, the prime-encoded quantum NLP engine, and the Prime-Embedded PageRank (PE-PageRank) algorithm into a unified quantum search engine. The system operates across three major layers: query processing, contextual analysis, and result ranking.

\section{Query Processing Layer}

\subsection{Hierarchical Node Parsing}
Kunferman's node-based parsing defines the structure of language using mutually exclusive categories. The parsing process can be represented as a directed graph:
\[
G = (V, E)
\]
where:
\begin{itemize}
    \item \( V \) represents nodes corresponding to words, phrases, or concepts.
    \item \( E \subseteq V \times V \) represents directed edges indicating hierarchical relationships.
\end{itemize}

The hierarchical node path for a query \(Q\) is generated using a traversal function:
\[
P(Q) = \{ v_1, v_2, \ldots, v_n \}, \quad v_i \in V
\]

\subsection{Prime-Encoding}
Each identified concept \(v_i\) is prime-encoded using a mapping function:
\[
\phi(v_i) = p_i, \quad p_i \text{ is a prime number}
\]

The prime-encoded vector for the query is:
\[
\Phi(Q) = [\phi(v_1), \phi(v_2), \ldots, \phi(v_n)]^T
\]

\section{Quantum Contextual Analysis Layer}

\subsection{Quantum State Representation}
The query is represented as a quantum state using a superposition of basis states:
\[
|\psi_Q\rangle = \sum_{i=1}^{n} \alpha_i |p_i\rangle
\]

where:
\begin{itemize}
    \item \( |p_i\rangle \) are basis states representing prime-encoded concepts.
    \item \( \alpha_i \in \mathbb{C} \) are complex coefficients representing probabilities.
\end{itemize}

\subsection{Probabilistic Context Resolution}
Using quantum gates, we apply probabilistic inference to evaluate context:
\[
U(\theta) |\psi_Q\rangle = \sum_{i=1}^{n} e^{i \theta_i} \alpha_i |p_i\rangle
\]

Measurement collapses the superposition to the most likely context:
\[
|\psi_{\text{context}}\rangle = |p_k\rangle \text{ with probability } |\alpha_k|^2
\]

\section{Prime-Embedded PageRank (PE-PageRank)}

\subsection{Prime-Weighted Transition Matrix}
The transition matrix \(M\) is defined using prime-weighted probabilities:
\[
M_{ij}(p) = \begin{cases}
\frac{1}{L(j) \cdot p(j)}, & \text{if there is a link from } j \text{ to } i \\
0, & \text{otherwise}
\end{cases}
\]

where:
\begin{itemize}
    \item \( L(j) \) is the number of outbound links from page \(j\).
    \item \( p(j) \) is the prime function associated with page \(j\).
\end{itemize}

\subsection{Prime-Enhanced Damping Factor}
The adaptive damping factor \(d_p(t)\) is given by:
\[
d_p(t) = d + \frac{1}{p(t)}
\]

where:
\begin{itemize}
    \item \(d\) is the classical damping factor (e.g., 0.85).
    \item \( p(t) \) is a prime number influenced by graph complexity.
\end{itemize}

\subsection{Prime-Embedded PageRank Calculation}
The PageRank vector \(R^{(t)}\) is updated using the iterative rule:
\[
R_i^{(t+1)} = \frac{1 - d_p(t)}{N} + d_p(t) \sum_{j \in L(i)} \frac{R_j^{(t)}}{L(j) \cdot p(j)}
\]

where:
\begin{itemize}
    \item \(N\) is the number of pages.
    \item \(L(i)\) is the set of pages linking to \(i\).
\end{itemize}

\subsection{Convergence Criteria}
The algorithm continues until the following convergence condition is satisfied:
\[
\| R^{(t+1)} - R^{(t)} \| < \epsilon_p
\]
where:
\[
\epsilon_p = \frac{1}{p(N)}
\]

\section{Final Ranking and Display}
After obtaining the PageRank vector \(R\), results are ranked using a hybrid score that integrates both contextual relevance and page importance:
\[
S_i = \lambda_1 C_i + \lambda_2 R_i
\]

where:
\begin{itemize}
    \item \(S_i\) is the final score of the \(i\)th result.
    \item \(C_i\) represents the context relevance derived from the quantum analysis.
    \item \(R_i\) is the PageRank score.
    \item \(\lambda_1, \lambda_2 \geq 0\) are adjustable parameters.
\end{itemize}

\section{Conclusion}
This framework provides a robust mathematical foundation for the quantum search engine, combining the interpretive strength of hierarchical parsing, the adaptive learning of prime-encoded tensors, and the efficient ranking of PE-PageRank. The integration of quantum inference further enables probabilistic context resolution, ensuring accurate and efficient search performance.

\nocite{*}
\bibliographystyle{plain}
\bibliography{references}

\end{document}